\subsection{Une approche en termes de sécurité globale}

Les crises majeures relèvent désormais d'un registre d'action multidimensionnel : militaire, civil, économique ou encore sociétal. Le comité s'est appuyé sur une vision en quatre cercles, proposés par le général de corps d'armée Benoît Durieux : défense militaire, défense nationale, sécurité nationale et sécurité internationale\footcite{IHEDN_perimetre_defense_nationale}. Cette grille d'analyse, illustrée sur la figure \ref{fig:cercles} en annexe \ref{annexe:figures_complementaires:cercles_structurants_action_IHEDN}, permet de structurer le champ d'action de l'IHEDN et d'articuler les enjeux de force armée, de souveraineté, de résilience nationale et de stabilité internationale.

Cette distinction permet de comprendre dans quelle mesure les crises contemporaines traversent plusieurs niveaux d'action à la fois. Une cyberattaque, une crise énergétique, une pandémie, une campagne de désinformation ou un événement climatique extrême peuvent produire des effets civils, économiques, institutionnels et stratégiques simultanés. Les frontières entre sphères civile et militaire deviennent alors plus poreuses, sans pour autant disparaître.

Ce cadre d'analyse évite de limiter la réflexion aux seuls dispositifs publics de gestion de crise, même si ceux-ci en constituent naturellement le socle. Il conduit à prendre en compte les ressources présentes dans des réseaux associatifs, professionnels et citoyens, dès lors que leur apport peut être identifié et organisé dans un cadre reconnu. À ce titre, l'\AssociationRegionaleIDF{} constitue un terrain d'observation intéressant : elle rassemble des membres issus d'administrations, des forces armées, d'entreprises, du monde académique et de la société civile, avec des expériences et des savoir-faire complémentaires.

\subsection{Menace, risque et crise : éléments de distinction}

La clarification des notions est nécessaire pour éviter de confondre des réalités proches, mais distinctes. Une menace désigne un facteur susceptible de porter atteinte à des intérêts, des personnes, des institutions ou des infrastructures. Elle peut être intentionnelle, lorsqu'elle procède d'une stratégie hostile, d'une action criminelle ou d'une pression politique, mais également résulter d'un aléa naturel, technologique ou sanitaire.

Le risque apparaît lorsque cette menace rencontre une vulnérabilité. Il dépend alors non seulement de la probabilité d'un événement, mais aussi du degré d'exposition des populations, de la fragilité des infrastructures, de la dépendance à certains systèmes ou encore de la capacité des organisations à anticiper et à absorber le choc. Cette approche met en évidence les interdépendances croissantes entre réseaux, territoires et institutions.

Dans cette perspective, la crise correspond moins au seul événement déclencheur qu'à la dégradation du fonctionnement ordinaire provoquée par la combinaison d'une menace, de vulnérabilités et d'effets en chaîne. Elle impose alors des arbitrages sous contrainte d'incertitude et nécessite une coordination exceptionnelle, souvent interinstitutionnelle. Cette distinction entre menace, risque et crise permet ainsi au comité de situer plus précisément la contribution possible de l'\ARSeize, principalement dans l'identification des vulnérabilités, la diffusion d'une culture de préparation et l'analyse des enseignements tirés des crises.

\subsection{L'évolution des menaces et l'hybridation des conflictualités}

Les crises contemporaines traduisent une évolution progressive des vulnérabilités stratégiques françaises.
Depuis les années 1980, les menaces auxquelles la France a été confrontée sont de natures diverses, allant de l'intensification des risques climatiques et énergétiques\footcite{AN:2022:Resilience_nationale} à la montée en puissance des cybermenaces\footcite{Berthier_hacktivisme} en passant par le développement des stratégies d'affrontement hybride\footcite{Gere_affrontement_hybride}.

Ces stratégies hybrides brouillent les repères habituels entre paix, crise et confrontation. Elles ne relèvent pas toujours d'une agression ouverte, mais peuvent installer une pression continue par la dissuasion, la dissimulation et la désinformation. La dissuasion vise à maintenir un rapport de force et de tension entre nations, sans franchir le seuil du conflit ouvert\footcite[p.89]{Gere_affrontement_hybride}. La dissimulation repose sur le recours à des acteurs indirects, tels que des mercenaires ou des groupes paraétatiques, ainsi que sur des actions clandestines comme les cyberattaques, afin de brouiller l'attribution des responsabilités\footcite[p.90]{Gere_affrontement_hybride}. Elle peut également prendre la forme d'incidents difficilement attribuables, tels que des intrusions de drones au-dessus de zones sensibles, observées récemment en Allemagne, en Pologne ou au Danemark\footcite{figaro_drones_europe}.

\clearpage

La désinformation occupe une place particulière dans cet ensemble, car elle agit sur la perception de la crise autant que sur la crise elle-même. Elle utilise les vecteurs de diffusion de l'information pour amplifier les succès du camp émetteur, minimiser ses échecs et fragiliser la perception et la cohésion du pays cible\footcite[p.91]{Gere_affrontement_hybride}. Les réseaux sociaux accélèrent cette dynamique : ils favorisent la circulation de récits hostiles, affaiblissent la confiance envers les institutions et peuvent amplifier les effets d'une crise déjà engagée\footcite{Sibony_reseaux_sociaux_et_guerre}. Dans un tel environnement, la résilience collective dépend aussi de la capacité des citoyens à comprendre l'information, à identifier les manipulations et à maintenir un lien de confiance avec les autorités légitimes.

Le tableau \ref{tab:vulnerabilites_strategiques_thematique} dresse un état des lieux thématiques des diverses crises auxquelles la France a été confrontée depuis 1980. Celui-ci témoigne de la pluralité des crises auxquelles l'État fut confronté et des vulnérabilités afférentes. Une analyse plus détaillée des crises mentionnées et de leurs effets est présentée dans le tableau \ref{tab:crises_majeures} en annexe \ref{annexe:tableau_crises_majeures}.

\begin{table}[H]
\centering
\StyledTableSetup

\stepcounter{footnote}\xdef\noteVulnerabilitesTerrorisme{\number\value{footnote}}
\stepcounter{footnote}\xdef\noteVulnerabilitesRisquesNaturels{\number\value{footnote}}
\stepcounter{footnote}\xdef\noteVulnerabilitesCrisesSanitaires{\number\value{footnote}}
\stepcounter{footnote}\xdef\noteVulnerabilitesInterdependances{\number\value{footnote}}

\begin{tabularx}{\linewidth}{@{}
p{2.8cm}
Y
Y
Y
Y
@{}}

\toprule
\textbf{Thématique} &
\textbf{Vulnérabilité dominante} &
\textbf{Type de menace ou de crise} &
\textbf{Effets sur la société} &
\textbf{Réponse développée} \\
\midrule

Terrorisme et sécurité intérieure\footnotemark[\noteVulnerabilitesTerrorisme] &
Fragilité des espaces publics et vulnérabilité des populations civiles &
Attentats terroristes et actions violentes organisées &
Sentiment d'insécurité, traumatisme collectif et tension sécuritaire &
Doctrine antiterroriste, renforcement du renseignement et dispositifs Vigipirate \\
\addlinespace

Risques naturels et infrastructures critiques\footnotemark[\noteVulnerabilitesRisquesNaturels] &
Dépendance aux réseaux énergétiques, logistiques et territoriaux &
Tempêtes majeures, accidents industriels et catastrophes naturelles &
Désorganisation territoriale et interruption des services essentiels &
Modernisation de la sécurité civile et politiques de prévention des risques \\
\addlinespace

Crises sanitaires et sociales\footnotemark[\noteVulnerabilitesCrisesSanitaires] &
Fragilité des systèmes sanitaires et tensions de cohésion sociale &
Canicule, pandémie, émeutes urbaines et mouvements sociaux &
Isolement, défiance institutionnelle et crise de cohésion &
Plans sanitaires, dispositifs de prévention et politiques de gestion de crise \\
\addlinespace

Cybermenaces et affrontements hybrides &
Vulnérabilité informationnelle et numérique &
Cyberattaques, désinformation et actions hybrides &
Défiance, diffusion des infox et brouillage de la perception des crises &
Approche de sécurité globale et renforcement des capacités cyber \\
\addlinespace

Interdépendances systémiques\footnotemark[\noteVulnerabilitesInterdependances] &
Imbrication des vulnérabilités sanitaires, énergétiques, économiques et géopolitiques &
Pandémie mondiale, guerre hybride et tensions énergétiques &
Crise multidimensionnelle affectant l'ensemble des fonctions sociales &
Développement d'une stratégie de résilience nationale et territoriale \\

\bottomrule
\end{tabularx}

\caption{Principales vulnérabilités stratégiques contemporaines et évolution des réponses publiques françaises.}
\label{tab:vulnerabilites_strategiques_thematique}
\end{table}

\footnotetext[\noteVulnerabilitesTerrorisme]{Attentats terroristes commis sur le territoire français depuis les années 1980, notamment les vagues d'attentats de 1986, 1995 et 2015.}
\footnotetext[\noteVulnerabilitesRisquesNaturels]{Tempêtes Lothar et Martin de décembre 1999, tempête Xynthia en 2010 et autres événements majeurs ayant affecté les infrastructures critiques.}
\footnotetext[\noteVulnerabilitesCrisesSanitaires]{Canicule de 2003, pandémie de COVID-19, émeutes urbaines de 2005 et mouvements sociaux récents.}
\footnotetext[\noteVulnerabilitesInterdependances]{Pandémie de COVID-19, guerre en Ukraine et tensions énergétiques associées.}

\clearpage

Dans un contexte de crises à la fois complexe et varié, garantir une gestion optimale et une diffusion de l'information appropriée est une nécessité. Afin de répondre à ces exigences, un cadre normatif de gestion de crises existe et permet d'apporter des solutions concrètes et pérennes.

\subsection{Une gestion de crise fortement structurée}

Sur le plan étatique, la gestion de crises se divise en quatre niveaux de gestion : politique, stratégique, tactique et opérationnel\footcite{niveaux_gestion_de_crise}.

Une telle organisation, détaillée dans le tableau \ref{tab:niveaux_crise}, permet de garantir une définition précise des responsabilités propres à chaque entité tout en coordonnant efficacement les acteurs engagés.

\begin{table}[H]
\centering
\small
\setlength{\tabcolsep}{4pt}
\resizebox{\linewidth}{!}{%
\begin{tabular}{@{}>{\raggedright\arraybackslash}p{2.5cm}>{\raggedright\arraybackslash}p{4.8cm}>{\raggedright\arraybackslash}p{4.8cm}>{\raggedright\arraybackslash}p{4.8cm}@{}}
\toprule
\textbf{Niveau} & \textbf{Définition simplifiée} & \textbf{Objectif principal} & \textbf{Exemples de profils} \\
\midrule

Politique &
Représente publiquement l'organisation et assume les décisions les plus sensibles. &
Maintenir la confiance, arbitrer les grandes décisions, protéger la réputation. &
Maires, ministres, PDG, présidents d'institution, top management, élus. \\

\addlinespace

Stratégique &
Pilote la réponse globale à la crise avec une vision à moyen et long terme. &
Assurer la continuité d'activité, préserver les actifs critiques, anticiper les impacts durables. &
Préfecture, direction générale, DGS, direction de site, comité de direction. \\

\addlinespace

Tactique &
Coordonne les moyens et met en œuvre les décisions stratégiques avec un appui technique. &
Adapter l'action à la situation, organiser les ressources, faire remonter les informations terrain. &
Commandants de secours, responsables HSE, référents sécurité, chefs d'équipe. \\

\addlinespace

Opérationnel &
Réalise les actions concrètes et immédiates pour répondre à la crise. &
Protéger les personnes et les biens, limiter les dommages immédiats. &
Pompiers, techniciens, agents de maintenance, forces de l'ordre, secouristes. \\

\bottomrule
\end{tabular}
}
\caption[Niveaux de coordination et de décision dans la gestion de crise.]{Niveaux de coordination et de décision dans la gestion de crise\footnotemark.}
\label{tab:niveaux_crise}
\end{table}

\footnotetext{Tableau issu de : \cite{niveaux_gestion_de_crise}.}

À cela s'ajoutent des dispositifs institutionnels de planification et de coordination tels que le dispositif \gls{orsec}\footcite{ORSEC_def}, les plans communaux de sauvegarde (\gls{pcs}) et les plans intercommunaux de sauvegarde (\gls{pics})\footcite{PCS_PICS_def}. Ces mécanismes, décrits plus précisément sur la figure \ref{fig:chaine_operationnelle} en annexe \ref{annexe:figures_complementaires:calendrier_questionnaire}, structurent l'anticipation, la préparation et la réponse aux crises à différentes échelles territoriales, tout en organisant la coordination entre autorités publiques, services de secours et acteurs concernés.

Ce constat conduit à préciser le positionnement de l'\ARSeize. 